% Part: methods
% Chapter: induction
% Section: induction-on-N

\documentclass[../../../include/open-logic-section]{subfiles}

\begin{document}

\olfileid{mth}{ind}{inN}

\olsection{$\Nat$పై ఆగమనం}

సరళమైన రూపంలో ఆగమనం అన్ని సహజ సంఖ్యల గురించి ఫలితాలను నిరూపించే పద్ధతి. $0$ నుంచి మొదలుపెట్టి~$1$ను పదేపదే కలుపుతూ వెళ్తే చివరికి ప్రతి సహజ సంఖ్యకు చేరతామనే విషయాన్ని ఉపయోగిస్తుంది. కాబట్టి ప్రతి సంఖ్యకు ఒక విషయం సత్యమని నిరూపించడానికి (1) అది $0$కు సత్యమని స్థాపించి, (2) ఒక సంఖ్య~$n$కు సత్యమైతే తరువాతి సంఖ్య~$n+1$కూ సత్యమని చూపవచ్చు. ``సంఖ్య~$n$కు $P$ లక్షణం ఉంది'' అన్నదాన్ని $P(n)$గా (అలాగే ``సంఖ్య~$k$కు $P$ లక్షణం ఉంది'' అన్నదాన్ని $P(k)$గా) సంక్షిప్తం చేస్తే, ప్రతి $n \in \Nat$కు $P(n)$ అని ఆగమన నిరూపణలో ఇవి ఉంటాయి:
\begin{enumerate}
\item $P(0)$ నిరూపణ; మరియు
\item ఏ~$k$కైనా $P(k)$ అయితే $P(k+1)$ అనే నిరూపణ.
\end{enumerate}
ఇది పూర్తిగా స్పష్టమయ్యేలా (1),~(2) రెండూ ఉన్నాయనుకుందాం. (1) ప్రకారం $P(0)$ సత్యం. (2) కూడా ఉంటే, ప్రత్యేకంగా $P(0)$ అయితే $P(0+1)$, అంటే $P(1)$ అని తెలుసు. ``ఏ~$k$కైనా $P(k)$ అయితే $P(k+1)$'' అనే సాధారణ వాక్యంలో~$k$ స్థానంలో~$0$ను పెట్టితే ఇది వస్తుంది. కాబట్టి మోడస్ పోనెన్స్ ప్రకారం~$P(1)$. మళ్లీ (2)లో ఇప్పుడు $k$ స్థానంలో $1$ పెట్టితే: $P(1)$ అయితే~$P(2)$. ఇప్పుడే~$P(1)$ను స్థాపించాం కనుక మోడస్ పోనెన్స్ ద్వారా~$P(2)$. ఇలా కొనసాగుతుంది. ఏ సంఖ్య~$n$కైనా ఈ ప్రక్రియను $n$~సార్లు చేస్తే చివరికి~$P(n)$కు చేరతాం. కనుక (1),~(2) కలిసి ఏ $n
\in \Nat$కైనా~$P(n)$ను స్థాపిస్తాయి.
\sourcecorrection{OLTEMTHINDN-001}{మూలంలో రెండవ దశకు $1$ను $n$ స్థానంలో పెట్టమంది; సాధారణ ఆగమన దశలో సూచిక $k$. లక్ష్యంలో అదే దశలో $k$ స్థానంలో $1$ను సరిగా పెట్టాం.}

ఒక ఉదాహరణ చూద్దాం. $n$ పాచికలతో ఎన్ని వేర్వేరు మొత్తాలు వేయవచ్చో తెలుసుకోవాలనుకుందాం. విచిత్రంగా అనిపించినా $0$ పాచికలతో మొదలుపెడదాం. పాచికలే లేకపోతే ``వేయగల'' మొత్తం ఒక్కటే: చుక్కలు ఏవీ లేవు; వాటి మొత్తం~$0$. కాబట్టి వేర్వేరు సాధ్యమైన ఫలితాల సంఖ్య~$1$. ఒకే పాచిక ఉంటే, అంటే $n=1$ అయితే $1$ నుంచి~$6$ వరకు ఆరు విలువలు. రెండు పాచికలతో $2$ నుంచి $12$ వరకు ఏ మొత్తమైనా వేయవచ్చు; అవి $11$ అవకాశాలు. మూడు పాచికలతో $3$ నుంచి $18$ వరకు ఏ సంఖ్యైనా, అంటే $16$ అవకాశాలు. $1$, $6$, $11$, $16$: ఒక నమూనా కనిపిస్తోంది; సమాధానం $5n+1$ కావచ్చా? నిజానికి $5n+1$ గరిష్ఠ సాధ్యమైన సంఖ్యే: $n$ పాచికలతో అతి తక్కువ మొత్తం $n$ (అన్నీ $1$లు), అతి ఎక్కువ మొత్తం $6n$ (అన్నీ $6$లు); వీటి మధ్య $5n+1$ సంఖ్యలే ఉన్నాయి.

\begin{thm}
  $n$ పాచికలతో $n$ నుంచి $6n$ వరకు ఉన్న $5n+1$ సాధ్యమైన విలువలన్నింటినీ వేయవచ్చు.
\end{thm}

\begin{proof}
$P(n)$ అనే వాదన: ``$n$~పాచికలతో $n$ నుంచి~$6n$ వరకు ఏ సంఖ్యనైనా వేయవచ్చు.'' ఆగమనాన్ని ఉపయోగించడానికి ఇవి నిరూపిద్దాం:
\begin{enumerate}
\item \emph{ఆగమన ఆధారం} $P(1)$: ఒకే పాచికతో $1$ నుంచి $6$ వరకు ఏ సంఖ్యనైనా వేయవచ్చు.
\item \emph{ఆగమన దశ}: అన్ని $k$లకు $P(k)$ అయితే~$P(k+1)$.
\end{enumerate}
శూన్య పాచికల సందర్భం కూడా విడిగా సత్యమే: చుక్కల మొత్తం సున్నా ఒక్కటే. దీన్ని విడిగా చేర్చితే, ఒకటి నుంచి పైకి వెళ్లే ఆధారం, దశతో పాటు సిద్ధాంతంలోని శూన్య సందర్భమూ కప్పబడుతుంది.
\sourcecorrection{OLTEMTHINDN-002}{మూల సిద్ధాంతం శూన్య పాచికల సందర్భాన్నీ కలిగి ఉన్నా, నిరూపణ $P(1)$ ఆధారం నుంచే మొదలైంది. ముందరి వివరణలో ఉన్న శూన్య పాచికల ఒకే మొత్తం విషయాన్ని నిరూపణలో విడిగా చేర్చి అన్ని సహజ సంఖ్యల పరిధిని కప్పాం.}

(1)ను $6$ వైపుల పాచికను చూసి నిరూపించవచ్చు. దానికి ఆరు వైపులూ ఉన్నాయి; $1$ నుంచి~$6$ వరకు ప్రతి సంఖ్య ఒక వైపుపై కనిపిస్తుంది. కాబట్టి ఒక పాచికతో $1$ నుంచి~$6$ వరకు ఏ సంఖ్యనైనా వేయవచ్చు.

(2) కోసం సోపాధిక పూర్వపక్షం, అంటే $P(k)$ను తీసుకుందాం. ఈ పరికల్పనను \emph{ఆగమన పరికల్పన} అంటారు. దాన్ని ఉపయోగించి~$P(k+1)$ను నిరూపించాలి. కష్టమైన భాగం: $k+1$ పాచికల మొత్తం సాధ్యమైన విలువలను $k$~పాచికల మొత్తం సాధ్యమైన విలువలు, అదనపు $k+1$వ పాచిక విలువల ద్వారా చూడటం. ఆగమన పరికల్పనను వాడాలంటే ఇదే చేయాలి.

ఆగమన పరికల్పన ప్రకారం $k$~పాచికలతో $k$ నుంచి~$6k$ వరకు ఏ సంఖ్యనైనా పొందవచ్చు. $(k+1)$వ పాచికపై~$1$ వస్తే మొత్తానికి $1$ కలుస్తుంది. కనుక $k$~పాచికలతో ఒక మొత్తం వేసి, $(k+1)$వ పాచికపై~$1$ వేసి $k+1$ నుంచి $6k+1$ వరకు ఏ విలువనైనా పొందవచ్చు. ఇంకా మిగిలినవి $6k+2$ నుంచి $6k+6$ వరకు. $k$ పాచికలన్నింటిపై $6$s వేసి, $(k+1)$వ పాచికపై $2$ నుంచి $6$ మధ్య సంఖ్య వేసి ఇవి పొందవచ్చు. కాబట్టి $k+1$ పాచికలతో $k+1$ నుంచి $6(k+1)$ వరకు ఏ సంఖ్యనైనా వేయవచ్చు; అంటే ఆగమన పరికల్పన~$P(k)$తో~$P(k+1)$ను నిరూపించాం.
\end{proof}

సంఖ్యలు, సమితులు మొదలైన వస్తువుల శ్రేణిని స్వయంగా ``ఆగమనాత్మకంగా'' నిర్వచించినప్పుడు వాటి గురించి నిరూపించడానికి ఆగమనం తరచుగా ఉపయోగపడుతుంది: అంటే $(n+1)$వ వస్తువును $n$వ వస్తువుతో నిర్వచిస్తాం. ఉదాహరణకు~$n$ వరకు సహజ సంఖ్యల మొత్తం~$s_n$ను ఇలా నిర్వచించవచ్చు:
\begin{align*}
  s_0 & = 0\\
  s_{n+1} & = s_n + (n+1)
\end{align*}
ఈ నిర్వచనం ఇస్తుంది:
\begin{align*}
  s_0 & = 0,\\
  s_1 & = s_0 + 1 && = 1,\\
  s_2 & = s_1 + 2 && = 1 + 2 = 3\\
  s_3 & = s_2 + 3 && = 1 + 2 + 3 = 6, \text{ etc.}
\end{align*}
ఇప్పుడు $s_n = n(n+1)/2$ అని ఆగమనంతో నిరూపించవచ్చు.

\begin{prop}
  $s_n = n(n+1)/2$.
\end{prop}

\begin{proof}
  (1) $s_0 = 0\cdot(0 + 1)/2$ అని, (2) $s_k =
  k(k+1)/2$ అయితే $s_{k+1} = (k+1)(k+2)/2$ అని నిరూపించాలి. (1) స్పష్టమే. (2) కోసం ఆగమన పరికల్పన $s_k = k(k+1)/2$ను తీసుకుందాం. దాన్ని ఉపయోగించి $s_{k+1} = (k+1)(k+2)/2$ అని చూపాలి.

  $s_{k+1}$ ఏమిటి? నిర్వచనం ప్రకారం $s_{k+1} = s_k + (k+1)$. ఆగమన పరికల్పన ప్రకారం $s_k = k(k+1)/2$. దీన్ని ముందరి సమీకరణలో ప్రతిస్థాపించి, తరువాత భిన్నాలతో కొంత అంకగణితం చేస్తే చాలు:
  \begin{align*}
    s_{k+1} & = \frac{k(k+1)}{2} + (k+1) = {}\\
    & = \frac{k(k+1)}{2} + \frac{2(k+1)}{2} = {}\\
    & = \frac{k(k+1) + 2(k+1)}{2} = {}\\
    & = \frac{(k+2)(k+1)}{2}.
  \end{align*}
\end{proof}

ఇక్కడ ముఖ్యమైన పాఠం: ఆగమనాత్మకంగా నిర్వచించిన శ్రేణి $a_n$ గురించి నిరూపిస్తున్నప్పుడు ఆగమనమే సహజ మార్గం. అలా నిర్వచించని సందర్భంలోనూ (పాచికల ఫలితాల మాదిరిగా),~$k+1$ సందర్భాన్ని~$k$ సందర్భంతో ఏదో విధంగా అనుసంధానించగలిగితే ఆగమనాన్ని వాడవచ్చు.

\end{document}
